% test-001.tex -- comando \spnum
%
% Compilar:  lualatex-dev test-001.tex
% Validar:   show-pdf-tags --xml test-001.pdf | rnv-wrapp
%            verapdf --flavour ua2 --format text test-001.pdf
%
% Cada caso es un parrafo numerado, para poder referirse a el al
% escuchar el PDF. Los comentarios "doc:" indican la lectura que
% documenta el manual de usuario; en el resto, la lectura se revisa
% escuchando. Las entradas invalidas (que detienen la compilacion)
% no van aqui.
\DocumentMetadata
  {
    lang=es-CL,
    pdfversion=2.0,
    pdfstandard=ua-2,
    tagging=on,
    tagging-setup={math/setup={mathml-SE}},
  }
\documentclass{article}
\usepackage[spanish]{babel}
\usepackage{unicode-math}
\usepackage{spintent}
\usepackage{hyperref}
\hypersetup
  {
    pdfdisplaydoctitle = true,
    pdftitle           = {test-001: spnum},
  }
\pagestyle{empty}
\begin{document}

\section*{A. Formas del argumento}

% doc: «23»
1. Entero sin signo: $\spnum{23}$.

% doc: «más 23»
2. Entero con signo positivo: $\spnum{+23}$.

% doc: «menos 23»
3. Entero con signo negativo: $\spnum{-23}$.

% doc: «3 coma 14»
4. Decimal con coma: $\spnum{3,14}$.

% doc: «3 punto 14»
5. Decimal con punto: $\spnum{3.14}$.

% doc: «menos 3 coma 14»
6. Decimal negativo: $\spnum{-3,14}$.

% doc: «3 coma 4 periódico»
7. Periódico puro: $\spnum{3;4}$.

% doc: «3 coma 1 con 4 periódico»
8. Periódico mixto: $\spnum{3,1;4}$.

% doc: «0 coma 1 con 23 periódico»
9. Periódico de dos cifras: $\spnum{0,1;23}$.

% doc: «menos 3 coma 1 con 4 periódico»
10. Periódico mixto negativo: $\spnum{-3,1;4}$.

\section*{B. Notación científica}

% doc: «3 coma 14 por 10 al cubo»
11. Con e minúscula: $\spnum{3,14e3}$.

% doc: «3 coma 14 por 10 al cubo»
12. Con E mayúscula: $\spnum{3,14E3}$.

% doc: «3 coma 14 por 10 a la menos 3 potencia»
13. Exponente negativo: $\spnum{3,14e-3}$.

% doc: «3 coma 14 por 10 elevado a la positivo 3 potencia»
14. Exponente con signo positivo: $\spnum{3,14e+3}$.

% doc: «5 por 10 cuadrado»
15. Entero con exponente: $\spnum{5e2}$.

% doc: «menos 3 coma 14 por 10 al cubo»
16. Número negativo con exponente: $\spnum{-3,14e3}$.

\section*{C. Llave cifras-sep}

% doc: «100000»
17. Por defecto: $\spnum{100000}$.

% doc: «100000»
18. Con true explícito: $\spnum[cifras-sep=true]{100000}$.

% doc: «1234567 coma 891234»
19. Parte entera y parte decimal: $\spnum{1234567,891234}$.

% doc: «100000»
20. Con false y espacios del usuario: $\spnum[cifras-sep=false]{10\,00\,00}$.

% doc: «100000»
21. Con false y sin espacios: $\spnum[cifras-sep=false]{100000}$.

\section*{D. Llave read-sign}

% doc: «menos 23»
22. Terse con signo negativo: $\spnum[read-sign=terse]{-23}$.

% doc: «23 negativo»
23. Clear con signo negativo: $\spnum[read-sign=clear]{-23}$.

% doc: «más 23»
24. Terse con signo positivo: $\spnum[read-sign=terse]{+23}$.

% doc: «23 positivo»
25. Clear con signo positivo: $\spnum[read-sign=clear]{+23}$.

% doc: «3 coma 14 negativo»
26. Clear con decimal negativo: $\spnum[read-sign=clear]{-3,14}$.

\section*{E. Llave read-period-sep}

% doc: «3 coma 4 periódico»
27. Con coma: $\spnum[read-period-sep=coma]{3;4}$.

% doc: «3 punto 4 periódico»
28. Con punto: $\spnum[read-period-sep=punto]{3;4}$.

% doc: «3 coma 1 con 4 periódico»
29. Con punto y parte decimal finita: $\spnum[read-period-sep=punto]{3,1;4}$.

\section*{F. Llave notation-sym}

% doc: «3 coma 14 por 10 al cubo»
30. Con cdot: $\spnum[notation-sym=cdot]{3,14e3}$.

% doc: «3 coma 14 por 10 al cubo»
31. Con times: $\spnum[notation-sym=times]{3,14e3}$.

\section*{G. Combinaciones de llaves}

% doc: «3 coma 14 negativo por 10 al cubo»
32. Clear y times: $\spnum[read-sign=clear,notation-sym=times]{-3,14e3}$.

% doc: «3 punto 4 periódico negativo»
33. Clear y punto: $\spnum[read-sign=clear,read-period-sep=punto]{-3;4}$.

% doc: «100000 negativo»
34. Clear y cifras-sep false: $\spnum[read-sign=clear,cifras-sep=false]{-10\,00\,00}$.

\section*{H. Dentro de fórmulas}

% doc: «3 coma 14 más 2 coma 5 es igual a 5 coma 64»
35. Suma e igualdad: $\spnum{3,14} + \spnum{2,5} = \spnum{5,64}$.

% doc: «2 negativo es menor que más 3»
36. Llaves distintas en una misma fórmula: $\spnum[read-sign=clear]{-2} < \spnum{+3}$.

% doc: «3 coma 14 por 10 al cubo»
37. En fórmula destacada:
\[ \spnum{3,14e3} \]

% doc: «la fracción con numerador 1 coma 5 y denominador 2»
38. En una fracción: $\frac{\spnum{1,5}}{\spnum{2}}$.

\end{document}
